\documentclass[11pt]{article}
\usepackage[a4paper,margin=1in]{geometry}
\usepackage{amsmath,amssymb,amsthm}
\usepackage[T1]{fontenc}
\usepackage{lmodern}
\usepackage[hidelinks]{hyperref}
\setlength{\emergencystretch}{3em}

\theoremstyle{plain}
\newtheorem{theorem}{Theorem}
\newtheorem{lemma}[theorem]{Lemma}
\newtheorem{proposition}[theorem]{Proposition}
\newtheorem{corollary}[theorem]{Corollary}
\theoremstyle{remark}
\newtheorem{remark}[theorem]{Remark}

\title{A Counterexample to Conjecture 00000000044}
\author{%
  Feiyu\thanks{Independent researcher.}
  \and
  Claude (Opus 5)\thanks{Anthropic.}%
}
\date{2026-09-12}

\begin{document}
\maketitle

\begin{abstract}
Conjecture 00000000044 asserts that every positive rational is a finite sum of
reciprocals of distinct primes, and that a greedy algorithm over prime
denominators always terminates. Both fail. If $q=\sum_{p\in S}1/p$ over a
finite set $S$ of distinct primes, then every prime dividing the reduced
denominator of $q$ lies in $S$; consequently $1/n$ is representable only when
$n=1$ or $n$ is prime, and $1/4$ is the smallest counterexample. The refutation
is formalized in Lean~4 against Mathlib.
\end{abstract}

\section{Statement}

Conjecture 00000000044 asks two questions and asserts a positive answer to
each:
\begin{enumerate}
  \item[(i)] every positive rational number is a finite sum of reciprocals of
    distinct primes;
  \item[(ii)] the greedy algorithm restricted to prime denominators always
    terminates in finitely many steps.
\end{enumerate}
Write
\[
  \mathcal R=\Bigl\{\textstyle\sum_{p\in S}\tfrac1p\;:\;S\subseteq\mathbb P
  \text{ finite}\Bigr\}
\]
for the set of values representable as in~(i), where $\mathbb P$ is the set of
primes and the elements of $S$ are distinct by virtue of $S$ being a set. For
$q\in\mathbb Q$ let $\operatorname{den}(q)$ denote the denominator of $q$ in
lowest terms.

\section{Two obstructions}

\begin{proposition}\label{prop:dvd}
Let $S\subseteq\mathbb P$ be finite and $q=\sum_{p\in S}1/p$. Then
\[
  \operatorname{den}(q)\ \Bigm|\ \prod_{p\in S}p .
\]
\end{proposition}

\begin{proof}
Induction on $S$. For $S=\varnothing$ the sum is $0$, whose denominator is $1$.
For the step, let $a\in\mathbb P$ and $a\notin S$. The denominator of a sum
divides the product of the denominators, so
\[
  \operatorname{den}\Bigl(\tfrac1a+\sum_{p\in S}\tfrac1p\Bigr)
  \;\Bigm|\;
  \operatorname{den}\bigl(\tfrac1a\bigr)\cdot
  \operatorname{den}\Bigl(\sum_{p\in S}\tfrac1p\Bigr)
  \;=\;a\cdot\operatorname{den}\Bigl(\sum_{p\in S}\tfrac1p\Bigr),
\]
using $\operatorname{den}(1/a)=a$ for $a$ prime. By the inductive hypothesis
this divides $a\prod_{p\in S}p$, which is the product over $S\cup\{a\}$.
\end{proof}

\begin{corollary}\label{cor:sqfree}
If $q\in\mathcal R$ then $\operatorname{den}(q)$ is squarefree.
\end{corollary}

\begin{proof}
A product of distinct primes is squarefree, and a divisor of a squarefree
number is squarefree.
\end{proof}

Corollary~\ref{cor:sqfree} already refutes (i): $\operatorname{den}(1/4)=4$ is
not squarefree, so $1/4\notin\mathcal R$. The next observation is sharper and
shows that the failure is not isolated.

\begin{proposition}\label{prop:mem}
Let $S\subseteq\mathbb P$ be finite and $q=\sum_{p\in S}1/p$. Then every prime
dividing $\operatorname{den}(q)$ belongs to $S$. Consequently
\begin{equation}\label{eq:lb}
  q\;\ge\sum_{\substack{r\ \mathrm{prime}\\ r\,\mid\,\operatorname{den}(q)}}\frac1r .
\end{equation}
\end{proposition}

\begin{proof}
Let $r$ be a prime dividing $\operatorname{den}(q)$. By
Proposition~\ref{prop:dvd}, $r\mid\prod_{p\in S}p$, so $r$ divides some
$p\in S$; as both are prime, $r=p\in S$. Thus the set of prime divisors of
$\operatorname{den}(q)$ is a subset of $S$, and \eqref{eq:lb} follows because
all the omitted terms $1/p$ are positive.
\end{proof}

\section{The conjecture is false}

\begin{theorem}\label{thm:main}
For $n\ge2$, if $1/n\in\mathcal R$ then $n$ is prime. In particular
$1/n\notin\mathcal R$ for every composite $n$, and the smallest counterexample
to (i) is
\[
  q=\tfrac14 .
\]
\end{theorem}

\begin{proof}
Let $n\ge2$ be composite and suppose $1/n\in\mathcal R$. We have
$\operatorname{den}(1/n)=n$. Let $r$ be the least prime factor of $n$. Since
$n$ is composite, $r<n$. By \eqref{eq:lb},
\[
  \frac1n\;\ge\sum_{\substack{s\ \mathrm{prime}\\ s\,\mid\,n}}\frac1s\;\ge\;\frac1r\;>\;\frac1n,
\]
a contradiction. (For $n$ prime the same inequality reads $1/n\ge1/n$ and is no
obstruction; indeed $1/n=\sum_{p\in\{n\}}1/p$.)
\end{proof}

So (i) fails for $1/4,\ 1/6,\ 1/8,\ 1/9,\ 1/10,\ 1/12,\dots$ --- for every
composite $n$. Statement (i) of the conjecture is false.

\begin{corollary}
Statement (ii) is false as well.
\end{corollary}

\begin{proof}
A terminating run of a procedure that repeatedly subtracts reciprocals of
distinct primes from a positive rational and stops at $0$ exhibits the starting
value as a finite sum of reciprocals of distinct primes, that is, as an element
of $\mathcal R$. Since $1/4\notin\mathcal R$, no such procedure terminates on
$1/4$ --- whatever rule it uses to choose the next prime, greedy or otherwise.
\end{proof}

\begin{remark}
The two obstructions are necessary conditions only, and we do not determine
$\mathcal R$. Corollary~\ref{cor:sqfree} and Proposition~\ref{prop:mem} say
nothing about, for instance, whether $1\in\mathcal R$: the denominator $1$ is
squarefree and has no prime divisors, so \eqref{eq:lb} is vacuous there. A
repaired conjecture would at minimum have to restrict to rationals $q$ with
$\operatorname{den}(q)$ squarefree satisfying \eqref{eq:lb}.
\end{remark}

\section{Formal verification}

The accompanying Lean~4 project formalizes the refutation against Mathlib. The
conjecture is stated in the formalization, so the final theorem is a
formalization of its negation rather than of facts that merely entail it.

\begin{itemize}
  \item \texttt{IsPrimeReciprocalSum}, \texttt{ConjectureHolds} --- membership
    in $\mathcal R$, and statement (i).
  \item \texttt{den\_dvd\_prod} --- Proposition~\ref{prop:dvd}, by induction on
    the finite set using Mathlib's \texttt{Rat.add\_den\_dvd}.
  \item \texttt{prod\_primes\_squarefree},
    \texttt{squarefree\_den\_of\_isPrimeReciprocalSum} ---
    Corollary~\ref{cor:sqfree}, with \texttt{Squarefree} Mathlib's notion.
  \item \texttt{primeFactors\_den\_subset}, \texttt{sum\_primeFactors\_den\_le}
    --- Proposition~\ref{prop:mem} and the inequality \eqref{eq:lb}.
  \item \texttt{not\_isPrimeReciprocalSum\_one\_div\_composite} ---
    Theorem~\ref{thm:main}.
  \item \texttt{conjecture\_00000000044\_false} ---
    $\lnot$\,\texttt{ConjectureHolds}, by instantiating at $q=1/4$.
  \item \texttt{no\_prime\_greedy\_terminates\_at\_one\_div\_four} --- the
    corollary for (ii), in the rule-independent form stated above.
\end{itemize}

The project builds with \texttt{lake build} on
\texttt{leanprover/lean4:v4.33.1} with Mathlib \texttt{v4.33.1}. It contains no
\texttt{sorry}, no \texttt{admit} and no \texttt{native\_decide}, and
introduces no axioms of its own:
\texttt{\#print axioms conjecture\_00000000044\_false} reports exactly the
three standard axioms \texttt{propext}, \texttt{Classical.choice} and
\texttt{Quot.sound}.

\end{document}
